\documentclass{article}
\usepackage{../../math_template}

\author{Jonathan Kasongo}
\title{Noether Sheet 6}


\begin{document}
\maketitle
\centerline{
These were my solutions to the Noether Mentoring Scheme Sheet 6. Problem
statements are abridged.
}

\begin{problem}{Problem 1}
The fibonacci sequence is the sequence $1,1,2,3,5,8,13,21,\ldots$, where
each term after the first two terms is the sum of the previous two.
\vspace{0.2cm}

(a) Show that there are infinitely many Fibonacci numbers that are the
mean of two different Fibonacci numbers.
\vspace{0.2cm}

(b) Show that there are infinitely many Fibonacci numbers that are the mean
of three different Fibonacci numbers.
\end{problem}

\begin{solution}{Solution}
(a) Let $F_n$ denote the $n$-th Fibonacci number. Call a Fibonacci number
$\textit{special}$ if it can be expressed as the mean of two different
Fibonacci numbers.
After some exploration
one may find that $F_5 = \frac12 (F_3 + F_6)$ and similarly,
$F_{11} = \frac12 (F_9 + F_{12})$. These results motivate the following
claim, \vspace{0.2cm}

\textit{Claim 1: For positive integer $n$,
$$
F_{3n + 2} = \frac12 (F_{3n} + F_{3n+3}).
$$}

\textit{Proof:} The proof is as follows,
\[
\begin{aligned}
F_{3n} + F_{3n+3} &= 2F_{3n+2} \\
F_{3n} + (F_{3n+1} + F_{3n+2}) &= F_{3n+2} + F_{3n+2}\\
F_{3n} + F_{3n+1} &= F_{3n+2}
\end{aligned}
\]
which is true by definition of the Fibonacci numbers. $\square$
\vspace{0.2cm}

Thus since there are infinitely many $n \in \mathbb{N}$, using claim 1
we can show that infinitely many Fibonacci numbers are special. $\square$
\\

(b) Call a Fibonacci number \textit{quirky} if it can be expressed as the
mean of three different Fibonacci numbers. Simimlarly to (a) after a bit
of exploration we may find that
$F_6 = \frac13 (F_4 + F_6 + F_7)$ and
$F_{14} = \frac13 (F_{12} + F_{14} + F_{15})$. \vspace{0.2cm}

\textit{Claim 2: For positive integer $n$,
$$
F_{8n-2} = \frac13 (F_{8n-4} + F_{8n-2} + F_{8n-1}).
$$
}

\textit{Proof:} The proof is as follows,
\[
\begin{aligned}
F_{8n-4} + F_{8n-2} + F_{8n-1} &= 3F_{8n-2}\\
F_{8n-4} + F_{8n-1} &= 2F_{8n-2}\\
F_{8n-4} + (F_{8n-2} + F_{8n-3}) &= 2F_{8n-2}\\
F_{8n-4} + F_{8n-3} &= F_{8n-2}\\
\end{aligned}
\]
which is true due to the definition of the Fibonacci numbers. $\square$

Thus since there are infinitely many $n \in \mathbb{N}$, using claim 2
we can show that infinitely many Fibonacci numbers are quirky.
$\blacksquare$
\end{solution}
\newpage

\begin{problem}{Problem 2}
Solve the simultaneous equations
\[
\begin{aligned}
x + \sqrt{y} &= 13,\\
xy &= 90.\\
\end{aligned}
\]
\end{problem}

\begin{solution}{Solution}
Square the first equation to find that $y = (13 - x)^2$. Sub into the
second equation to get
\[
\begin{aligned}
x(13 - x)^2 &= 90 \\
x^3 - 26x^2 + 169x - 90 &= 0 \\
(x - 10)(x^2 - 16x + 9) &= 0 \\
\end{aligned}
\]

which solves to give $x = 10, 8 + \sqrt{55}, 8 - \sqrt{55}$. But now one can
simply check that the only solutions are
$(x, y) = (10, 9), (8 - \sqrt{55}, 80 + 10\sqrt{55})$. $\blacksquare$
\end{solution}
\newpage

\begin{problem}{Problem 3}
Let $ABC$ be a triangle.
\vspace{0.2cm}

(a) Prove that the perpendicular bisectors of the three sides meet at a
point. This is the \textit{circumcentre} of $\triangle ABC$, usually
denoted $O$. \vspace{0.2cm}

Prove also that the circle through $A, B$ and $C$ has it's centre at $O$.
\vspace{0.2cm}

(b) Prove that the angle bisectors of $\angle A, \angle B$ and $\angle C$
meet at a point. This is the \textit{incentre} of $\triangle ABC$, usually
denoted $I$. \vspace{0.2cm}

(c) Prove that the altitudes of the triangle meet at a point. This is the
\textit{orthocentre} of the $\triangle ABC$, usually denoted $H$.
\end{problem}

\begin{solution}{Solution}
(a) Construct points $L, M, N$ to be the midpoints of the sides
$BC, AB, AC$ respectively. Let the side lengths be $a,b,c$ is the usual
convention. Let $O$ be the intersection of the perpendicular bisectors
of $AC, AB$. Such a point must exist because two lines intersect if they
aren't parallel, since $AC$ and $AB$ aren't parallel, their perpendicular
bisectors aren't parallel either. We have the diagram,

\begin{center}
\begin{asy}
/* Geogebra to Asymptote conversion, documentation at artofproblemsolving.com/Wiki go to User:Azjps/geogebra */
import graph; size(210);
real labelscalefactor = 0.5; /* changes label-to-point distance */
pen dps = linewidth(0.7) + fontsize(10); defaultpen(dps); /* default pen style */
pen dotstyle = black; /* point style */
real xmin = -1.1271641961608405, xmax = 10.833480545268682, ymin = -0.20225179001776172, ymax = 7.880244504948259;  /* image dimensions */
pen qqwuqq = rgb(0.,0.39215686274509803,0.);
pair A = (1.,1.), B = (9.,1.), C = (3.,7.), M = (5.,1.), L = (6.,4.), O = (5.,3.), N = (2, 4);
pair MCL = (4.5, 5.5), MBL = (7.5, 2.5), MBM = (7, 1), MAM = (3, 1), MAN = (1.5, 2.5), MCN = (2.5, 5.5);

filldraw((5.255837953371849,1.)--(5.255837953371849,1.2558379533718482)--(5.,1.2558379533718482)--M--cycle, qqwuqq + opacity(0.10000000149011612),  qqwuqq);
filldraw((1.9190969355429004,3.7572908066287014)--(2.161806128914199,3.676387742171602)--(2.2427091933712986,3.9190969355429006)--(2.,4.)--cycle, qqwuqq + opacity(0.10000000149011612),  qqwuqq);
/* draw figures */
draw(A--B);
draw(C--B);
draw(C--A);
draw((2.,4.)--O,  linetype("4 4"));
draw(O--M,  linetype("4 4"));
draw(O--L,  linetype("4 4"));
draw(O--B);
draw(O--C);
draw(A--O);
/* dots and labels */
dot(A,dotstyle);
label("$A$", A, NW * labelscalefactor);
dot(B,dotstyle);
label("$B$", B, NE * labelscalefactor);
dot(C,dotstyle);
label("$C$", C, NE * labelscalefactor);
dot(M, dotstyle);
label("$M$", M, NW * labelscalefactor);
dot(L, dotstyle);
label("$L$", L, NE * labelscalefactor);
dot((2.,4.), dotstyle);
label("$N$", N, NE * labelscalefactor);
dot(O, dotstyle);
label("$O$", O, NE * labelscalefactor);

label("$c/2$", MBM, SE * labelscalefactor);
label("$c/2$", MAM, SE * labelscalefactor);
label("$b/2$", MAN, NW * labelscalefactor);
label("$b/2$", MCN, NW * labelscalefactor);
label("$a/2$", MCL, NE * labelscalefactor);
label("$a/2$", MBL, NE * labelscalefactor);

clip((xmin,ymin)--(xmin,ymax)--(xmax,ymax)--(xmax,ymin)--cycle);
/* end of picture */
\end{asy}
\end{center}

And our goal is to show that $\angle OLB = 90^\circ$.
\vspace{0.2cm}

\textit{Claim 1: If $OC = OB$ then this implies $\angle OLB = 90^\circ$.}
\vspace{0.2cm}

\textit{Proof:} Let $\angle OLB = \theta$. Due to cosine rule and
angles on a straight line summing to 180,
\[
\begin{aligned}
OC^2 &= \left( \frac{a}{2} \right)^2 + OL^2 -
2 \left(\frac{a}{2}\right) (OL) \cos (180 - \theta) \\
OB^2 &= \left( \frac{a}{2} \right)^2 + OL^2 -
2 \left(\frac{a}{2}\right) (OL) \cos (\theta) \\
\end{aligned}
\]
and so if $OB = OC$ then $\cos (\theta) = \cos(180 - \theta) =
-\cos(\theta)$. But this means that $\cos \theta = 0$ and since
$\theta \in (0, 180)$ we have $\theta = 90^\circ$ only. $\square$ \\

Denote $x = OB, y = OC, h = OM$.
By Pythagoras,
\[
\begin{aligned}
x^2 &= h^2 + \left( \frac{c}{2} \right)^2 \\
y^2 &= \left( \frac{b}{2} \right)^2 + ON^2 \\
&= \left( \frac{b}{2} \right)^2 + \left[
OA^2 - \left( \frac{b}{2} \right)^2 \right]  \\
&= OA^2 = \left( \frac{c}{2} \right)^2 + h^2 \\
\end{aligned}
\]

So this must mean that $x = y$ and due to claim 1 we are done. Note that
$OA^2 = (c/2)^2 + h^2$ also by Pythagoras and so $OA = OB = OC$. Since the
point $O$ is equidistant from all three points, the circle through $A, B, C$
has centre $O$ and radius $OA$. $\square$
\\

(b) Let $I$ intersection of the angle bisectors $B_A$ of $\angle A$ and
$B_B$ of $\angle B$. We need to show that such a point must exist.
\vspace{0.2cm}

WLOG assume the point configuration as seen in the diagram. If two
lines aren't parallel then they must intersect at a point. Suppose for
the sake of contradiction that those two bisectors are parallel.
Let $\theta$ be the angle $B_A$ makes with the side $AB$ (measuring
anti-clockwise), since $B_B$ is parallel it also makes an angle of $\theta$
of with the side $AB$ (measuring anti-clockwise). That means that
measuring clockwise the angle it makes is $180 - \theta$. But this angle
is exactly $\frac12 \angle B$ so $\angle B = 360 - 2\theta$, similarly
$\angle A = 2\theta$. But clearly this means that the interior angle sum
in $\triangle ABC$ is greater than 180, a contradiction. \vspace{0.2cm}

Let $P, N, M$ be the feet of the perpendiculars from $I$ to sides
$AC, BC, AB$. Our goal is to show that $\angle PCI = \angle NCI$.
\vspace{0.2cm}

\begin{center}
\begin{asy}
/* Geogebra to Asymptote conversion, documentation at artofproblemsolving.com/Wiki go to User:Azjps/geogebra */
import graph; size(210);
real labelscalefactor = 0.5; /* changes label-to-point distance */
pen dps = linewidth(0.7) + fontsize(10); defaultpen(dps); /* default pen style */
pen dotstyle = black; /* point style */
real xmin = -2.6827729217113703, xmax = 12.811002999016145, ymin = -1.604348280411282, ymax = 8.732561309184055;  /* image dimensions */
pen qqwuqq = rgb(0.,0.39215686274509803,0.);
pair A = (0.,0.), B = (9.,0.), C = (6.,7.), I = (5.301885675714489,2.4385223772067337), M = (5.301885675714488,0.), P = (3.450421460804758,4.025491704272217);

filldraw((5.301885675714488,0.45068546827615547)--(4.851200207438334,0.450685468276156)--(4.851200207438333,0.)--M--cycle, qqwuqq + opacity(0.10000000149011612),  qqwuqq);
filldraw((3.1571192459086723,3.6833057868934507)--(3.499305163287439,3.390003571997365)--(3.7926073781835243,3.7321894893761316)--P--cycle, qqwuqq + opacity(0.10000000149011612),  qqwuqq);
filldraw(arc(A,0.7648386019207299,360.,384.69935267749776)--(0.,0.)--cycle, qqwuqq + opacity(0.10000000149011612),  qqwuqq);
filldraw(arc(A,0.9560482524009124,384.69935267749776,409.3987053549955)--(0.,0.)--cycle, qqwuqq + opacity(0.10000000149011612),  qqwuqq);
filldraw(arc(B,0.7648386019207299,146.59929525682412,180.)--(9.,0.)--cycle, qqwuqq + opacity(0.10000000149011612),  qqwuqq);
filldraw(arc(B,0.8923117022408517,473.1985905136481,506.59929525682406)--(9.,0.)--cycle, qqwuqq + opacity(0.10000000149011612),  qqwuqq);
filldraw((7.330200955060078,3.896197771526483)--(6.833106585139211,3.683157327274683)--(7.046147029391011,3.1860629573538155)--(7.543241399311878,3.3991034016056156)--cycle, qqwuqq + opacity(0.10000000149011612),  qqwuqq);
/* draw figures */
draw(A--C);
draw(C--B);
draw(B--A);
draw(A--I,  linetype("4 4"));
draw(I--B,  linetype("4 4"));
draw(I--C,  linetype("4 4"));
draw(I--M);
draw(I--P);
draw(arc(A,0.7648386019207299,360.,384.69935267749776),  qqwuqq);
draw(arc(A,0.6819810867126509,360.,384.69935267749776),  qqwuqq);
draw(arc(A,0.9560482524009124,384.69935267749776,409.3987053549955),  qqwuqq);
draw(arc(A,0.8731907371928334,384.69935267749776,409.3987053549955),  qqwuqq);
draw(arc(B,0.7648386019207299,146.59929525682412,180.),  qqwuqq);
draw(arc(B,0.6819810867126509,146.59929525682412,180.),  qqwuqq);
draw(arc(B,0.5991235715045719,146.59929525682412,180.),  qqwuqq);
draw(arc(B,0.8923117022408517,473.1985905136481,506.59929525682406),  qqwuqq);
draw(arc(B,0.8094541870327726,473.1985905136481,506.59929525682406),  qqwuqq);
draw(arc(B,0.7265966718246936,473.1985905136481,506.59929525682406),  qqwuqq);
draw(I--(7.543241399311878,3.3991034016056156));
/* dots and labels */
dot(A, dotstyle);
label("$A$", (0.05731495367236327,0.0963241562213689), NW * labelscalefactor);
dot(B,dotstyle);
label("$B$", (9.056915836272951,0.12181877628539327), NE * labelscalefactor);
dot(C,dotstyle);
label("$C$", (6.048550668718081,7.132839293892092), NE * labelscalefactor);
dot(I, dotstyle);
label("$I$", (5.347448616957411,2.5438076823677074), NE * labelscalefactor);
dot(M, dotstyle);
label("$M$", (5.347448616957411,0.0963241562213689), NE * labelscalefactor);
dot(P, dotstyle);
label("$P$", (3.4990886623156476,4.124474126337217), NW * labelscalefactor);
dot((7.543241399311878,3.3991034016056156), dotstyle);
label("$N$", (7.590975182591553,3.499855934768621), NE * labelscalefactor);
clip((xmin,ymin)--(xmin,ymax)--(xmax,ymax)--(xmax,ymin)--cycle);
/* end of picture */
\end{asy}
\end{center}

\textit{Claim 1: A point $A$ lies on the bisector of $PQ$ and $PR$ if and
only if the perpendiculars from $A$ to $PQ, PR$ are the same length.}
\vspace{0.2cm}

\textit{Proof:} Let (1) be the statement
"A point $A$ lies on the bisector of $PQ$ and $PR$".
Let (2) be the statement "The perpendiculars from $A$ to $PQ, PR$ are the
same length". \vspace{0.2cm}

First we show that $(1) \implies (2)$.
Let $F_Q, F_R$ be the feet of the perpendiculars from
$A$ to $PQ, PR$ respectively. Since $\angle APQ = \angle APR$ by assumption
and $\angle P F_Q A = \angle P F_R A = 90^\circ$, and angles in a triangle
sum
to 180 means that their third angles are the same also, the triangles
$\triangle APF_Q$ and $\triangle APF_R$ are similar by AAA criterion.
In fact since both triangles share the hypotenuse $AP$, the two triangles
are congruent. This gives $AF_Q = AF_R$.
\vspace{0.2cm}

Next we show that $(2) \implies (1)$. Define $F_Q, F_R$ as we did before.
Since $\angle A F_Q P = \angle A F_R P = 90^\circ$, $AF_R = AF_Q$ by
assumption and the triangles $\triangle AF_Q P$ and $\triangle AF_R P$
both share hypotenuse $AP$, the two triangles are congruent by SAS
criterion. But this means that $\angle APF_R = \angle APF_Q$. $\square$
\\

Now using the fact that two equal angles in two triangles implies that their
third angles are equal, the triangles $\triangle PIA$ and
$\triangle MIA$ are similar and so are the triangles
$\triangle MIB$ and $\triangle NIB$. Moreover since each pair of triangles
shares a hypotenuse the each pair of two triangles are actually congruent.
This allows us to conclude $IM = PI = NI$. \vspace{0.2cm}

Now we have $NI = PI$ so claim 2 tells us that $I$ lies on the angle
bisector of $(PC, CN)$. That means that $CI$ bisects $\angle PCN$ and
so $\angle PCI = \angle NCI$ as desired. $\square$
\\

(c) Let $H$ be the intersection of the altitudes through $C$ and through
$B$. Two lines intersect if they aren't parallel. The altitude
through $C$ is perpendicular to $AB$, and the altitude through $B$ is
perpendicular $AC$, so if the two altitudes were parallel then this would
imply that $AC \parallel AB$ which is impossible for a non-degenerate
triangle.
\vspace{0.2cm}

Let $F_B, F_C$ be the feet of the altitudes through $B,C$
respectively. Let $F_A$ be the intersection of the line through $(A, H)$
and side $BC$. Our goal is to show $\angle AF_A B = 90^\circ$.

\begin{center}
\begin{asy}
/* Geogebra to Asymptote conversion, documentation at artofproblemsolving.com/Wiki go to User:Azjps/geogebra */
import graph; size(230);
real labelscalefactor = 0.5; /* changes label-to-point distance */
pen dps = linewidth(0.7) + fontsize(10); defaultpen(dps); /* default pen style */
pen dotstyle = black; /* point style */
real xmin = -1.3619020684590244, xmax = 14.729499480406975, ymin = -0.7953404609680161, ymax = 10.383421362698996;  /* image dimensions */
pen qqwuqq = rgb(0.,0.39215686274509803,0.);
pair A = (1.,1.), B = (10.,1.), C = (4.,9.), H = (4.,3.25), F_A = (6.76,5.32), F_B = (2.1095890410958904,3.958904109589041), F_C = (4.,1.);

filldraw((4.575357853277581,1.)--(4.575357853277581,1.575357853277581)--(4.,1.575357853277581)--F_C--cycle, qqwuqq + opacity(0.10000000149011612),  qqwuqq);
filldraw((2.6483133866631787,3.7568824800013076)--(2.850335016250912,4.295606825568596)--(2.3116106706836237,4.497628455156329)--F_B--cycle, qqwuqq + opacity(0.10000000149011612),  qqwuqq);
/* draw figures */
draw(A--B);
draw(C--B);
draw(C--A);
draw(A--F_A);
draw(B--F_B);
draw(F_C--C);
draw(circle((2.5,2.125), 1.875));
draw(F_B--F_C,  linetype("4 4"));
draw(circle((7.,5.), 5.));
/* dots and labels */
dot(A,dotstyle);
label("$A$", (1.0631228384646503,1.1641542877834476), NE * labelscalefactor);
dot(B,dotstyle);
label("$B$", (10.062411243777303,1.1641542877834476), NE * labelscalefactor);
dot(C,dotstyle);
label("$C$", (4.057461114373381,9.17075446032199), NW * labelscalefactor);
dot(H, dotstyle);
label("$H$", (4.057461114373381,3.3773608395420696), NE * labelscalefactor);
dot(F_A, dotstyle);
label("$F_A$", (6.823969304071665,5.444105193022547), NE * labelscalefactor);
dot(F_B, dotstyle);
label("$F_B$", (2.169726114343964,4.093398253346329), NW * labelscalefactor);
dot(F_C, dotstyle);
label("$F_C$", (4.057461114373381,1.1316071326105268), SE * labelscalefactor);
clip((xmin,ymin)--(xmin,ymax)--(xmax,ymax)--(xmax,ymin)--cycle);
/* end of picture */
\end{asy}
\end{center}

\textit{Claim 1: $A F_B H F_C$ is a cyclic quadilateral.}
\vspace{0.2cm}

\textit{Proof:} We have $\angle AF_C H = \angle AF_B H = 90^\circ$, so
these two opposite angles sum to 180. Now since angles in a quadrilateral
sum to 360 this forces $\angle F_B A F_C + \angle F_B H F_C = 180^\circ$,
which means that $A F_B H F_C$ is cyclic. $\square$ \\

\textit{Claim 2: $C F_B F_C B$ is a cylic quadrilateral.}
\vspace{0.2cm}

\textit{Proof:} Firstly notice that $\angle ACF_C = 90 - \angle A$.
Then since $A F_B H F_C$ is cyclic, we know that
$\angle F_B H F_C = 180 - \angle A$ which means that
$\angle F_C H B = \angle A$ due to angles on a straight line sum to 180.
Then we get that $\angle H B F_C = 90 - \angle A$. \vspace{0.2cm}

Now the key step is to use the fact that $\angle F_B B F_C = \angle
F_B C F_C$, and then by the converse of "Angles subtending the same segment
are equal" it must be the case that $F_B, F_C, B, C$ all lie on the same
circle. $\square$ \\

Now since $C F_B F_C B$ is cyclic, we know that
$\angle F_B F_C B = 180 - \angle C$. Since $\angle HF_CB$ is right-angled
we deduce $\angle F_B F_C H = 90 - \angle C$. Now consider the cyclic
quadrilateral $A F_B H F_C$.
Since angles subtending
the same segment are equal we have $\angle F_B A H = 90 - \angle C$.
Finally we can look at $\triangle C F_A A$ and deduce that
$\angle C F_A A = 180 - (\angle C + 90 - \angle C) = 90^\circ$.
as desired. $\blacksquare$
\end{solution}
\newpage

\begin{problem}{Problem 4}
The five tetrominoes are shown below. They can be turned either way up.

\begin{center}
\begin{asy}
/* Geogebra to Asymptote conversion, documentation at artofproblemsolving.com/Wiki go to User:Azjps/geogebra */
import graph; size(170);
real labelscalefactor = 0.5; /* changes label-to-point distance */
pen dps = linewidth(0.7) + fontsize(10); defaultpen(dps); /* default pen style */
pen dotstyle = black; /* point style */
real xmin = -0.8282552137749006, xmax = 15.645367724743695, ymin = -0.6384843590483772, ymax = 10.613025335560076;  /* image dimensions */

/* draw figures */
draw((2.,7.)--(4.,7.));
draw((4.,7.)--(4.,5.));
draw((4.,5.)--(2.,5.));
draw((2.,7.)--(2.,5.));
draw((3.,7.)--(3.,5.),  linetype("4 4"));
draw((2.,6.)--(4.,6.),  linetype("4 4"));
draw((6.,7.)--(8.,7.));
draw((8.,7.)--(8.,6.));
draw((7.,6.)--(8.,6.));
draw((7.,6.)--(7.,5.));
draw((7.,5.)--(5.,5.));
draw((5.,5.)--(5.,6.));
draw((5.,6.)--(6.,6.));
draw((6.,6.)--(6.,7.));
draw((6.,6.)--(7.,6.),  linetype("4 4"));
draw((6.,6.)--(6.,5.),  linetype("4 4"));
draw((7.,7.)--(7.,6.),  linetype("4 4"));
draw((9.,7.)--(9.,5.));
draw((9.,5.)--(12.,5.));
draw((12.,5.)--(12.,6.));
draw((12.,6.)--(10.,6.));
draw((10.,6.)--(10.,7.));
draw((10.,7.)--(9.,7.));
draw((9.,6.)--(10.,6.),  linetype("4 4"));
draw((10.,6.)--(10.,5.),  linetype("4 4"));
draw((11.,6.)--(11.,5.),  linetype("4 4"));
draw((3.,3.)--(7.,3.));
draw((7.,3.)--(7.,2.));
draw((7.,2.)--(3.,2.));
draw((3.,2.)--(3.,3.));
draw((4.,3.)--(4.,2.),  linetype("4 4"));
draw((5.,2.)--(5.,3.),  linetype("4 4"));
draw((6.,3.)--(6.,2.),  linetype("4 4"));
draw((8.,2.)--(11.,2.));
draw((11.,3.)--(11.,2.));
draw((8.,3.)--(8.,2.));
draw((9.,3.)--(8.,3.));
draw((10.,3.)--(11.,3.));
draw((9.,4.)--(10.,4.));
draw((10.,4.)--(10.,3.));
draw((9.,4.)--(9.,3.));
draw((9.,3.)--(9.,2.),  linetype("4 4"));
draw((10.,3.)--(10.,2.),  linetype("4 4"));
draw((9.,3.)--(10.,3.),  linetype("4 4"));
/* dots and labels */
clip((xmin,ymin)--(xmin,ymax)--(xmax,ymax)--(xmax,ymin)--cycle);
/* end of picture */
\end{asy}
\end{center}

(a) Show that a $4 \times 6$ and a $3\times 8$ rectangle can be tiled by
six tetrominoes, five of them different. Which tetromino have you used
twice? \vspace{0.2cm}

(b) Show that no rectangle can be tiled by the set of five each used once.
\end{problem}

\begin{solution}{Solution}
(a) Here are the two rectangles.

\begin{center}
\begin{asy}
/* Geogebra to Asymptote conversion, documentation at artofproblemsolving.com/Wiki go to User:Azjps/geogebra */
import graph; size(170);
real labelscalefactor = 0.5; /* changes label-to-point distance */
pen dps = linewidth(0.7) + fontsize(10); defaultpen(dps); /* default pen style */
pen dotstyle = black; /* point style */
real xmin = 0.0825866388662369, xmax = 16.78750866076997, ymin = 0.5331405008056004, ymax = 9.350176266212774;  /* image dimensions */
pen ttttff = rgb(0.2,0.2,1.); pen dcrutc = rgb(0.8627450980392157,0.0784313725490196,0.23529411764705882); pen qqwuqq = rgb(0.,0.39215686274509803,0.); pen ffwwqq = rgb(1.,0.4,0.); pen ffdxqq = rgb(1.,0.8431372549019608,0.); pen qqqqcc = rgb(0.,0.,0.8); pen qqccqq = rgb(0.,0.8,0.); pen wwccff = rgb(0.4,0.8,1.);

filldraw((2.,1.)--(2.,2.)--(3.,2.)--(3.,3.)--(4.,3.)--(4.,2.)--(5.,2.)--(5.,1.)--cycle, ttttff + opacity(0.3499999940395355),  ttttff);
filldraw((3.,5.)--(4.,5.)--(4.,4.)--(5.,4.)--(5.,2.)--(4.,2.)--(4.,3.)--(3.,3.)--cycle, dcrutc + opacity(0.3499999940395355),  dcrutc);
filldraw((2.,6.)--(3.,6.)--(3.,2.)--(2.,2.)--cycle, qqwuqq + opacity(0.3499999940395355),  qqwuqq);
filldraw((2.,8.)--(4.,8.)--(4.,6.)--(2.,6.)--cycle, black + opacity(0.3499999940395355));
filldraw((2.,8.)--(2.,9.)--(5.,9.)--(5.,7.)--(4.,7.)--(4.,8.)--cycle, ffwwqq + opacity(0.3499999940395355),  ffwwqq);
filldraw((4.,7.)--(5.,7.)--(5.,4.)--(4.,4.)--(4.,5.)--(3.,5.)--(3.,6.)--(4.,6.)--cycle, ffdxqq + opacity(0.3499999940395355),  ffdxqq);
filldraw((7.,1.)--(7.,3.)--(9.,3.)--(9.,1.)--cycle, red + opacity(0.3499999940395355),  red);
filldraw((8.,3.)--(8.,4.)--(10.,4.)--(10.,1.)--(9.,1.)--(9.,3.)--cycle, qqqqcc + opacity(0.3499999940395355),  qqqqcc);
filldraw((10.,1.)--(11.,1.)--(11.,5.)--(10.,5.)--cycle, qqccqq + opacity(0.3499999940395355),  qqccqq);
filldraw((7.,3.)--(7.,6.)--(8.,6.)--(8.,5.)--(9.,5.)--(9.,4.)--(8.,4.)--(8.,3.)--cycle, ffdxqq + opacity(0.3499999940395355),  ffdxqq);
filldraw((9.,4.)--(9.,6.)--(10.,6.)--(10.,7.)--(11.,7.)--(11.,5.)--(10.,5.)--(10.,4.)--cycle, black + opacity(0.3499999940395355));
filldraw((7.,6.)--(7.,7.)--(10.,7.)--(10.,6.)--(9.,6.)--(9.,5.)--(8.,5.)--(8.,6.)--cycle, wwccff + opacity(0.3499999940395355),  wwccff);
/* draw figures */
draw((2.,1.)--(2.,2.),  ttttff);
draw((2.,2.)--(3.,2.),  ttttff);
draw((3.,2.)--(3.,3.),  ttttff);
draw((3.,3.)--(4.,3.),  ttttff);
draw((4.,3.)--(4.,2.),  ttttff);
draw((4.,2.)--(5.,2.),  ttttff);
draw((5.,2.)--(5.,1.),  ttttff);
draw((5.,1.)--(2.,1.),  ttttff);
draw((3.,5.)--(4.,5.),  dcrutc);
draw((4.,5.)--(4.,4.),  dcrutc);
draw((4.,4.)--(5.,4.),  dcrutc);
draw((5.,4.)--(5.,2.),  dcrutc);
draw((5.,2.)--(4.,2.),  dcrutc);
draw((4.,2.)--(4.,3.),  dcrutc);
draw((4.,3.)--(3.,3.),  dcrutc);
draw((3.,3.)--(3.,5.),  dcrutc);
draw((2.,6.)--(3.,6.),  qqwuqq);
draw((3.,6.)--(3.,2.),  qqwuqq);
draw((3.,2.)--(2.,2.),  qqwuqq);
draw((2.,2.)--(2.,6.),  qqwuqq);
draw((2.,8.)--(4.,8.));
draw((4.,8.)--(4.,6.));
draw((4.,6.)--(2.,6.));
draw((2.,6.)--(2.,8.));
draw((2.,8.)--(2.,9.),  ffwwqq);
draw((2.,9.)--(5.,9.),  ffwwqq);
draw((5.,9.)--(5.,7.),  ffwwqq);
draw((5.,7.)--(4.,7.),  ffwwqq);
draw((4.,7.)--(4.,8.),  ffwwqq);
draw((4.,8.)--(2.,8.),  ffwwqq);
draw((4.,7.)--(5.,7.),  ffdxqq);
draw((5.,7.)--(5.,4.),  ffdxqq);
draw((5.,4.)--(4.,4.),  ffdxqq);
draw((4.,4.)--(4.,5.),  ffdxqq);
draw((4.,5.)--(3.,5.),  ffdxqq);
draw((3.,5.)--(3.,6.),  ffdxqq);
draw((3.,6.)--(4.,6.),  ffdxqq);
draw((4.,6.)--(4.,7.),  ffdxqq);
draw((7.,1.)--(7.,3.),  red);
draw((7.,3.)--(9.,3.),  red);
draw((9.,3.)--(9.,1.),  red);
draw((9.,1.)--(7.,1.),  red);
draw((8.,3.)--(8.,4.),  qqqqcc);
draw((8.,4.)--(10.,4.),  qqqqcc);
draw((10.,4.)--(10.,1.),  qqqqcc);
draw((10.,1.)--(9.,1.),  qqqqcc);
draw((9.,1.)--(9.,3.),  qqqqcc);
draw((9.,3.)--(8.,3.),  qqqqcc);
draw((10.,1.)--(11.,1.),  qqccqq);
draw((11.,1.)--(11.,5.),  qqccqq);
draw((11.,5.)--(10.,5.),  qqccqq);
draw((10.,5.)--(10.,1.),  qqccqq);
draw((7.,3.)--(7.,6.),  ffdxqq);
draw((7.,6.)--(8.,6.),  ffdxqq);
draw((8.,6.)--(8.,5.),  ffdxqq);
draw((8.,5.)--(9.,5.),  ffdxqq);
draw((9.,5.)--(9.,4.),  ffdxqq);
draw((9.,4.)--(8.,4.),  ffdxqq);
draw((8.,4.)--(8.,3.),  ffdxqq);
draw((8.,3.)--(7.,3.),  ffdxqq);
draw((9.,4.)--(9.,6.));
draw((9.,6.)--(10.,6.));
draw((10.,6.)--(10.,7.));
draw((10.,7.)--(11.,7.));
draw((11.,7.)--(11.,5.));
draw((11.,5.)--(10.,5.));
draw((10.,5.)--(10.,4.));
draw((10.,4.)--(9.,4.));
draw((7.,6.)--(7.,7.),  wwccff);
draw((7.,7.)--(10.,7.),  wwccff);
draw((10.,7.)--(10.,6.),  wwccff);
draw((10.,6.)--(9.,6.),  wwccff);
draw((9.,6.)--(9.,5.),  wwccff);
draw((9.,5.)--(8.,5.),  wwccff);
draw((8.,5.)--(8.,6.),  wwccff);
draw((8.,6.)--(7.,6.),  wwccff);
/* dots and labels */
clip((xmin,ymin)--(xmin,ymax)--(xmax,ymax)--(xmax,ymin)--cycle);
/* end of picture */
\end{asy}
\end{center}

In each case we used the "T" tetromino twice. $\square$
\\

(b) Observe the following diagram.

\begin{center}
\begin{asy}
/* Geogebra to Asymptote conversion, documentation at artofproblemsolving.com/Wiki go to User:Azjps/geogebra */
import graph; size(150);
real labelscalefactor = 0.5; /* changes label-to-point distance */
pen dps = linewidth(0.7) + fontsize(10); defaultpen(dps); /* default pen style */
pen dotstyle = black; /* point style */
real xmin = 0.5071096746123463, xmax = 10.65398055532054, ymin = 0.4805785164044289, ymax = 8.357208349815396;  /* image dimensions */
pen ududff = rgb(0.30196078431372547,0.30196078431372547,1.);

filldraw((1.,5.)--(2.,5.)--(2.,4.)--(1.,4.)--cycle, red + opacity(0.3499999940395355),  red);
filldraw((1.,4.)--(2.,4.)--(2.,3.)--(1.,3.)--cycle, ududff + opacity(0.3499999940395355),  ududff);
filldraw((1.,3.)--(2.,3.)--(2.,2.)--(1.,2.)--cycle, red + opacity(0.3499999940395355),  red);
filldraw((1.,2.)--(2.,2.)--(2.,1.)--(1.,1.)--cycle, ududff + opacity(0.3499999940395355),  ududff);
filldraw((3.,3.)--(3.,2.)--(4.,2.)--(4.,3.)--cycle, ududff + opacity(0.3499999940395355),  ududff);
filldraw((4.,4.)--(5.,4.)--(5.,3.)--(4.,3.)--cycle, ududff + opacity(0.3499999940395355),  ududff);
filldraw((4.,2.)--(4.,3.)--(5.,3.)--(5.,2.)--cycle, red + opacity(0.3499999940395355),  red);
filldraw((5.,4.)--(6.,4.)--(6.,3.)--(5.,3.)--cycle, red + opacity(0.3499999940395355),  red);
filldraw((8.,3.)--(8.,4.)--(9.,4.)--(9.,3.)--cycle, red + opacity(0.3499999940395355),  red);
filldraw((7.,2.)--(7.,3.)--(8.,3.)--(8.,2.)--cycle, red + opacity(0.3499999940395355),  red);
filldraw((9.,2.)--(9.,3.)--(10.,3.)--(10.,2.)--cycle, red + opacity(0.3499999940395355),  red);
filldraw((8.,2.)--(8.,3.)--(9.,3.)--(9.,2.)--cycle, ududff + opacity(0.3499999940395355),  ududff);
filldraw((2.,8.)--(3.,8.)--(3.,7.)--(2.,7.)--cycle, red + opacity(0.3499999940395355),  red);
filldraw((3.,7.)--(4.,7.)--(4.,6.)--(3.,6.)--cycle, red + opacity(0.3499999940395355),  red);
filldraw((3.,7.)--(3.,8.)--(4.,8.)--(4.,7.)--cycle, ududff + opacity(0.3499999940395355),  ududff);
filldraw((2.,6.)--(2.,7.)--(3.,7.)--(3.,6.)--cycle, ududff + opacity(0.3499999940395355),  ududff);
filldraw((6.,7.)--(7.,7.)--(7.,6.)--(6.,6.)--cycle, ududff + opacity(0.3499999940395355),  ududff);
filldraw((8.,7.)--(9.,7.)--(9.,6.)--(8.,6.)--cycle, ududff + opacity(0.3499999940395355),  ududff);
filldraw((7.,7.)--(8.,7.)--(8.,6.)--(7.,6.)--cycle, red + opacity(0.3499999940395355),  red);
filldraw((8.,7.)--(8.,8.)--(9.,8.)--(9.,7.)--cycle, red + opacity(0.3499999940395355),  red);
/* draw figures */
draw((1.,5.)--(2.,5.),  red);
draw((2.,5.)--(2.,4.),  red);
draw((2.,4.)--(1.,4.),  red);
draw((1.,4.)--(1.,5.),  red);
draw((1.,4.)--(2.,4.),  ududff);
draw((2.,4.)--(2.,3.),  ududff);
draw((2.,3.)--(1.,3.),  ududff);
draw((1.,3.)--(1.,4.),  ududff);
draw((1.,3.)--(2.,3.),  red);
draw((2.,3.)--(2.,2.),  red);
draw((2.,2.)--(1.,2.),  red);
draw((1.,2.)--(1.,3.),  red);
draw((1.,2.)--(2.,2.),  ududff);
draw((2.,2.)--(2.,1.),  ududff);
draw((2.,1.)--(1.,1.),  ududff);
draw((1.,1.)--(1.,2.),  ududff);
draw((3.,3.)--(3.,2.),  ududff);
draw((3.,2.)--(4.,2.),  ududff);
draw((4.,2.)--(4.,3.),  ududff);
draw((4.,3.)--(3.,3.),  ududff);
draw((4.,4.)--(5.,4.),  ududff);
draw((5.,4.)--(5.,3.),  ududff);
draw((5.,3.)--(4.,3.),  ududff);
draw((4.,3.)--(4.,4.),  ududff);
draw((4.,2.)--(4.,3.),  red);
draw((4.,3.)--(5.,3.),  red);
draw((5.,3.)--(5.,2.),  red);
draw((5.,2.)--(4.,2.),  red);
draw((5.,4.)--(6.,4.),  red);
draw((6.,4.)--(6.,3.),  red);
draw((6.,3.)--(5.,3.),  red);
draw((5.,3.)--(5.,4.),  red);
draw((8.,3.)--(8.,4.),  red);
draw((8.,4.)--(9.,4.),  red);
draw((9.,4.)--(9.,3.),  red);
draw((9.,3.)--(8.,3.),  red);
draw((7.,2.)--(7.,3.),  red);
draw((7.,3.)--(8.,3.),  red);
draw((8.,3.)--(8.,2.),  red);
draw((8.,2.)--(7.,2.),  red);
draw((9.,2.)--(9.,3.),  red);
draw((9.,3.)--(10.,3.),  red);
draw((10.,3.)--(10.,2.),  red);
draw((10.,2.)--(9.,2.),  red);
draw((8.,2.)--(8.,3.),  ududff);
draw((8.,3.)--(9.,3.),  ududff);
draw((9.,3.)--(9.,2.),  ududff);
draw((9.,2.)--(8.,2.),  ududff);
draw((2.,8.)--(3.,8.),  red);
draw((3.,8.)--(3.,7.),  red);
draw((3.,7.)--(2.,7.),  red);
draw((2.,7.)--(2.,8.),  red);
draw((3.,7.)--(4.,7.),  red);
draw((4.,7.)--(4.,6.),  red);
draw((4.,6.)--(3.,6.),  red);
draw((3.,6.)--(3.,7.),  red);
draw((3.,7.)--(3.,8.),  ududff);
draw((3.,8.)--(4.,8.),  ududff);
draw((4.,8.)--(4.,7.),  ududff);
draw((4.,7.)--(3.,7.),  ududff);
draw((2.,6.)--(2.,7.),  ududff);
draw((2.,7.)--(3.,7.),  ududff);
draw((3.,7.)--(3.,6.),  ududff);
draw((3.,6.)--(2.,6.),  ududff);
draw((6.,7.)--(7.,7.),  ududff);
draw((7.,7.)--(7.,6.),  ududff);
draw((7.,6.)--(6.,6.),  ududff);
draw((6.,6.)--(6.,7.),  ududff);
draw((8.,7.)--(9.,7.),  ududff);
draw((9.,7.)--(9.,6.),  ududff);
draw((9.,6.)--(8.,6.),  ududff);
draw((8.,6.)--(8.,7.),  ududff);
draw((7.,7.)--(8.,7.),  red);
draw((8.,7.)--(8.,6.),  red);
draw((8.,6.)--(7.,6.),  red);
draw((7.,6.)--(7.,7.),  red);
draw((8.,7.)--(8.,8.),  red);
draw((8.,8.)--(9.,8.),  red);
draw((9.,8.)--(9.,7.),  red);
draw((9.,7.)--(8.,7.),  red);
/* dots and labels */
clip((xmin,ymin)--(xmin,ymax)--(xmax,ymax)--(xmax,ymin)--cycle);
/* end of picture */
\end{asy}
\end{center}

Thus if we were to tile the tetrominoes on a chessboard, then every
tetromino except the "T" tetromino will cover 2 white squares and 2 black
squares. The "T" tetromino covers 3 white squares and 1 black square
or 3 black squares and one white square. In any case if we tile all five
tetrominoes together on a rectangular chessboard we will cover 2 more white
squares than black squares (or 2 more black squares than white squares).
However this is impossible, since a chessboard has an equal number of
white and black squares. Therefore it is not possible to tile a rectangle
using each tetromino exactly once. $\blacksquare$
\end{solution}
\newpage

\begin{problem}{Problem 5}
(a) Suppose the numbers $1,2,\ldots, n$ are to be split into three
non-empty subsets. How many ways are there to achieve this?
\vspace{0.2cm}

(b) What happens if we insist that each subset consists of a run of at least
one consecutive number?
\end{problem}

\begin{solution}{Solution}
(a) Clearly $n\geq3$.
Suppose the first subset has size $a$ the second has size $b$ and the
third has size $c$.
\begin{itemize}
\item Since $b,c \geq 1$ we find that $a$ is at most $n-2$.
\item Since $c\geq 1$ we have that $a+b \leq n-1$ and so $b$ is at most
$n-1-a$.
\item Finally there is exactly one choice for $c$ that is $c = n-a-b$.
\end{itemize}
Putting this together, the number of choices turns out to be
$$
S = \sum_{a=1}^{n-2}\left[ \binom{n}{a} \times \sum_{b=1}^{n-a-1}
\binom{n-a}{b} \times 1 \right].
$$

By Binomial theorem,
$$
\sum_{b=0}^{n-a} \binom{n-a}{b} = (1+1)^{n-a} = 2^{n-a}
$$
and so we can re-write $S$ as
$$
S = \sum_{a=1}^{n-2} \left[ \binom{n}{a} (2^{n-a} - 2)\right].
$$
We apply a similar observation
$$
\sum_{a=0}^n \binom{n}{a} 2^{-a} = \left(1 + \frac12\right)^n =
\left(\frac32\right)^n
\quad
\text{and}
\quad
\sum_{a=0}^n \binom{n}{a} = (1+1)^n = 2^n.
$$
Thus
\[
\begin{aligned}
S &= 2^n \left[ \left(\frac32\right)^n - \binom{n}{0} -
\binom{n}{n-1}\frac{1}{2^{n-1}} - \binom{n}{n} \frac{1}{2^n} \right] -
2\left[2^n - \binom{n}{0} - \binom{n}{n-1} - \binom{n}{n}\right]\\
&= \left[3^n - 2^n - 2n - 1 \right] - 2\left[2^n - n - 2\right]\\
&= 3^n - 3(2^n) + 3
\end{aligned}
\]
and so there are exactly $3^n - 3(2^n) + 3$ ways to split the numbers.
$\square$\\

(b) To split the numbers in this way we can think about placing two
"dividers" in the spots underlined below

$$
\text{
\Large
$1\ \_ \ 2\ \_ \ 3\ \_ \ 4 \ \_ \ 5 \ \_ \ldots \_ \ (n-1) \ \_ \ n$
}
$$
There are $(n-1)$ places to put the dividers so the answer is
$\binom{n-1}{2}$. $\blacksquare$
\end{solution}
\newpage

\begin{problem}{Problem 6}
(a) In a group of six academics, each pair corresponds only about one topic. Moreover, the only two
topics are geometry and algebra. Prove that there must be three academics such that each pair of
those three correspond with each other about the same topic.
\vspace{0.2cm}

(b) In a group of seventeen academics, each pair corresponds only about one topic. Moreover, the only
three topics are geometry, algebra and number theory. Prove that there must be three academics
such that each pair of those three correspond with each other about the same topic.
\end{problem}

\begin{solution}{Solution}
(a)
\textit{Claim 1: Consider the complete graph with 6 vertices, $K_6$. Colour
every edge either red or blue. Then there must exist a monochromatic
3-cycle in $K_6$.}
\vspace{0.2cm}

\textit{Proof:} We proceed by way of contradiction.
On $K_6$ every vertex has degree 5. Consider vertex $i$.
Suppose there are $r$ red edges and $b$ blue edges that connect to $i$.
At least one of $r,b$ must be $\geq 3$ because otherwise
$$
\deg i = 5 = r+b \leq 2 + 2 = 4
$$
which is impossible. Suppose that the edges from vertex $i$ to
vertices $a,b,c$ are all the same colour, say WLOG they are all red. Since
there doesn't exist a monochromatic 3-cycle on $K_6$ all of the edges
$a - b$, $b - c$ and $c - a$ must be blue. But now notice that
there is a monochromatic 3-cycle between $a - b - c$ where all edges are
blue. Contradiction. $\square$
\\

If we model each academic as a vertex on a graph, colour an edge red
if a pair corresponds about geometry and colour an edge blue if a pair
corresponds about algebra, then the problem reduces to showing that claim 1
is true, as we have done. $\square$
\\

(b) \textit{Claim 2: Consider the complete graph with 17 vertices,
$K_{17}$. Colour every edge either red, green or blue. Then there must
exist a monochromatic 3-cycle in $K_{17}$}
\vspace{0.2cm}

\textit{Proof:} We proceed by way of contradiction. On $K_{17}$ every vertex
has degree 16. Consider a vertex $i$. Suppose that there are
$r$ red edges, $b$ blue edges and $g$ green edges that connect to $i$.
At least one of $r,b,g$ must be $\geq 6$ because otherwise
$$
\deg i = 16 = r + b + g \leq 5 + 5 + 5 = 15
$$
which is impossible. Suppose that the edges from vertex $i$ to the vertices
$a,b,c,d,e,f$ are all the same colour, say WLOG they are red. We note that
any edges between two of these vertices cannot be red, because there are
no 3-cycles on $K_{17}$. \vspace{0.2cm}

Thus the vertices $a,b,c,d,e,f$ actually form a 2-coloured complete
sub-graph with 6 vertices, since edge colours can only be blue or green.
However due to claim 1, there must exist a monochromatic 3-cycle on this
sub-graph, contradiction. $\square$\\

If we model each academic as a vertex on a graph, colour an edge red
if a pair corresponds about geometry and colour an edge blue if a pair
corresponds about algebra and colour an edge green if a pair corresponds
about number theory, then the problem reduces to showing that claim 2
is true, as we have done. $\blacksquare$
\\

\textit{Comments: This is really the first time I've done any graph theory
in maths, though previously I've come across graphs in computer science
and through solving problems on atcoder.jp and codeforces.com etc. Actually
today I was recommended a youtube video that was a lecture on Ramsey
theory, and
I ended up finding out that part (a) is equivalent to showing that
$R(3,3) = 6$ and part (b) is equivalent to showing that $R(3,3,3)=17$.}
\end{solution}
\newpage

\begin{problem}{Problem 7}
Let $G$ be a convex quadrilateral. Show that there is a point $X$ in the
plane of $G$ with the property that every straight line through $X$ divides
$G$ into two regions of equal area if and only if $G$ is a parallelogram.
\end{problem}

\begin{solution}{Partial Solution}
First we show the following statement to be true. If $G$ is a parallelogram
then there exists some point $X$, with the property that any line through
$X$ splits $G$ into two regions of equal area. Let $G$ have vertices
$ABCD$.
\vspace{0.2cm}

Take $X = AD \cap BC$ like so.
\begin{center}
\begin{asy}
/* Geogebra to Asymptote conversion, documentation at artofproblemsolving.com/Wiki go to User:Azjps/geogebra */
import graph; size(210);
real labelscalefactor = 0.5; /* changes label-to-point distance */
pen dps = linewidth(0.7) + fontsize(10); defaultpen(dps); /* default pen style */
pen dotstyle = black; /* point style */
real xmin = -1.071933600271589, xmax = 8.5653337946408, ymin = -1.0230073202821973, ymax = 4.907180676820011;  /* image dimensions */
pen ttttff = rgb(0.2,0.2,1.);
pair A = (0.,0.), B = (6.,0.), C = (2.,4.), D = (8.,4.), X = (4.,2.);

filldraw(C--D--B--A--cycle, ttttff + opacity(0.10000000149011612),  ttttff);
/* draw figures */
draw(C--D,  ttttff);
draw(D--B,  ttttff);
draw(B--A,  ttttff);
draw(A--C,  ttttff);
draw(C--B);
draw(A--D);
draw((xmin, -4.965414931958289*xmin + 21.861659727833157)--(xmax, -4.965414931958289*xmax + 21.861659727833157),  linetype("4 4")); /* line */
/* dots and labels */
dot(A, dotstyle);
label("$A$", (0.05096696358781762,0.10455017641659202), NE * labelscalefactor);
dot(B,dotstyle);
label("$B$", (6.055239796089695,0.1292591181141306), NE * labelscalefactor);
dot(C,dotstyle);
label("$C$", (2.0523912410884435,4.119753202266612), NE * labelscalefactor);
dot(D,dotstyle);
label("$D$", (8.044309602741551,4.119753202266612), NE * labelscalefactor);
dot(X, dotstyle);
label("$X$", (4.05381551858907,2.093619983068448), NE * labelscalefactor);
clip((xmin,ymin)--(xmin,ymax)--(xmax,ymax)--(xmax,ymin)--cycle);
/* end of picture */
\end{asy}
\end{center}


\textit{Claim 1: $BC$ and $AD$ are lines of symmetry of $G$.}
\vspace{0.2cm}

\textit{Proof:} We show that $\triangle AXB \equiv \triangle CXD$.
\begin{itemize}
\item $AB = CD$ by properties of a parallelogram.
\item $\angle BAX = \angle CDX$ because of alternate angles.
\item $\angle ABX = \angle XCD$ because of alternate angles.
\end{itemize}
Thus the two triangles are indeed congruent by the ASA criterion. This
allows us to conclude that $AX = DX$. Now since $A,X,D$ are collinear,
$D$ is in fact the reflection of the point $A$ across $BC$. Moreover
points $C, B$ are invariant under the transformation of reflection across
$BC$, since they lie on this line. This means that $G$ is in fact unchanged
under this reflection and so $BC$ is a line of symmetry
of $G$. By a similar arguement we can show that $AD$ is also a line of
symmetry of $G$.
$\square$
\\

Now consider some line $l$ passing through $X$.
\vspace{0.2cm}

\textit{Case 1: The line $l$ doesn't pass through any of $A, B, C$ or $D$.}
We assume WLOG that $l$ passes through sides $AB$ and $CD$. We have
already seen that $AX = DX$ and $CX = BX$. The points $A, D$ lie on
opposite sides of the line $l$, and the same is true for the points
$B, C$. Now since $A, X, D$ are collinear and $l$ passes though $X$,
$D$ is the reflection of $A$ across $l$.
Similarly $C$ is the reflection of $B$ across $l$. \vspace{0.2cm}

Suppose that $l$ cuts through $CD$ and $AB$ at points $H, K$ respectively.
If we reflect $ACHK$ across $l$ it maps onto $BDHK$. Since
reflection transformations preserve areas, we have that the area of
$ACHK$ and the area of $BDHK$ are equal.
\\

\textit{Case 2: The line $l$ passes through two diagonally opposite
vertices.} Then $l$ either passes through $BC$ or $AD$. Due to claim 1,
we established that these are lines of symmetry of $G$, so the two regions
formed map onto each other upon reflection across $l$. Again since
reflections preserve area, the area of the two regions are equal.
\\

Next we show the following statement to be true. For some convex
quadrilateral $G$, if there exists a point $X$ with the property that
any line through $X$ divides $G$ into two regions of equal area then $G$
must be a parallelogram. Call such a point $X$ \textit{special}.
\vspace{0.2cm}

We instead prove the contrapostive of this statement, which is logically
equivalent. We show that if $G$ is not a parallelogram, then there does not
exist a special point $X$ in the plane.\vspace{0.2cm}

Let us suppose that $G$ has vertices $ABCD$.

\begin{center}
\begin{asy}
/* Geogebra to Asymptote conversion, documentation at artofproblemsolving.com/Wiki go to User:Azjps/geogebra */
import graph; size(170);
real labelscalefactor = 0.5; /* changes label-to-point distance */
pen dps = linewidth(0.7) + fontsize(10); defaultpen(dps); /* default pen style */
pen dotstyle = black; /* point style */
real xmin = -2.0991317996570658, xmax = 10.954488410387674, ymin = -3.312073815761595, ymax = 5.71415039149418;  /* image dimensions */
pen ttttff = rgb(0.2,0.2,1.); pen qqzzqq = rgb(0.,0.6,0.);
pair A = (0.,0.), B = (4.9498390358302675,0.8517969150759968), C = (6.2173683456576105,2.317931116762521), D = (2.,4.);

filldraw(A--D--C--B--cycle, ttttff + opacity(0.10000000149011612),  ttttff);
//filldraw(arc(D,0.5059542717071605,-116.56505117707799,-21.74425128726117)--(2.,4.)--cycle, qqzzqq + opacity(0.10000000149011612),  qqzzqq);
//filldraw(arc(B,0.5059542717071605,49.155326516450366,189.76415485058544)--(4.9498390358302675,0.8517969150759968)--cycle, red + opacity(0.10000000149011612),  red);
/* draw figures */
draw(A--D,  ttttff);
draw(D--C,  ttttff);
draw(C--B,  ttttff);
draw(B--A,  ttttff);
/* dots and labels */
dot(A, dotstyle);
label("$A$", (0.03599522694715149,0.07781980467639213), NE * labelscalefactor);
dot(B,dotstyle);
label("$B$", (4.9943470896773245,0.9480611520127112), SE * labelscalefactor);
dot(C,dotstyle);
label("$C$", (6.259232768945225,2.415328539963482), NE * labelscalefactor);
dot(D,dotstyle);
label("$D$", (2.0395741429075067,4.105215807465404), NE * labelscalefactor);
//label("$\beta$", (2.0699313992099366,3.6802142192313876), SW * labelscalefactor,qqzzqq);
//label("$\alpha$", (4.771727210126174,1.0694901772224301), NW * labelscalefactor,red);
clip((xmin,ymin)--(xmin,ymax)--(xmax,ymax)--(xmax,ymin)--cycle);
/* end of picture */
\end{asy}
\end{center}

\textit{Claim 1: For any line $l$ splitting $G$ into two regions of non-zero
area we can translate $l$ onto precisely one line $l'$ such that $l'$
splits $G$ into two regions of equal areas.}
\vspace{0.2cm}

\textit{Proof: } We may rotate the entire configuration so that $l$ is
parallel to the $y$-axis, because rotations leave lengths and areas
unchanged. Now $l$ splits $G$ into a left region with area $L$ and a right
region with area $R$.
Suppose the parallelogram has area 1. If we move $l$ to be as rightmost
as possible then $L \rightarrow 1$ and $R \rightarrow 0$, similarly
if we move $l$ to be as leftmost as possible then $L \rightarrow 0$ and
$R \rightarrow 1$. Recalling that $L + R = 1$, by applying
intermediate value theorem coupled with the fact that moving the line in
one direction changes $L, R$ monotonically
there must be some unique point at which $L = R = 1/2$.
$\square$ \\

If I take some line $M$ that splits $G$ into two regions
of equal area, then for any point $P$ that doesn't lie on $M$, we can take
the line through $P$ that is parallel to $M$. Then since $P \neq M$ due to
claim 1, this line doesn't split $G$ into two regions of equal area.
Thus any point $P$ not lying on $M$ isn't special.
\vspace{0.2cm}

We may make the exact same statement but this time replacing the line $M$
with some line $T$ that isn't parallel to $M$. So any point $P$ not lying
on $T$ isn't special either. The only point remaining that could be special
is the point $Y = T \cap M$, since it is the only point lying on both
$T$ and $M$. \vspace{0.2cm}

The solution is nearly finished, we just need to patch up this line
passing through $Y$ case.
\end{solution}
\newpage


\begin{problem}{Problem 8}
Let $p \equiv 3 \ (\text{mod } 4)$ be a prime.\vspace{0.2cm}

(a) Suppose that $x^4 \equiv 1 \ (\text{mod } p)$. Prove that
$x^2 \equiv 1 \ (\text{mod } p)$.\vspace{0.2cm}

(b) Therefore show that there are no solutions to
$x^2 + 1 \equiv 0 \ (\text{mod } p)$. \vspace{0.2cm}

(c) Prove that there are an infinite number of primes
$\equiv 1 \ (\text{mod } 4)$.
\end{problem}

\begin{solution}{Solution}
(a) Write $p = 4k+3$ for $k \in \mathbb{Z}_{\geq0}$.
We are given $x^4 \equiv 1 \ (\text{mod } p)$.
Note that since $x^4$ leaves a non-zero residue modulo $p$, the numbers
$x$ and $p$ must be coprime. Thus by Fermat's little theorem
\[
\begin{aligned}
x^{4k+2} &\equiv 1 \ (\text{mod } p) \\
(x^4)^k \cdot x^2 &\equiv 1 \ (\text{mod } p) \\
1 \cdot x^2 &\equiv 1 \ (\text{mod } p) \\
\end{aligned}
\]
and so $x^2 \equiv 1 \ (\text{mod } p)$ as desired. $\square$\\

(b) \textit{Solution 1.} Suppose for the sake of contradiction that
$x^2 \equiv -1 \ (\text{mod } p)$. We find that

\[
\begin{aligned}
x^2 \cdot x^2 &\equiv (-1) \cdot (-1) \ (\text{mod } p) \\
x^4 &\equiv 1 \ (\text{mod } p)
\end{aligned}
\]

However in part (a) we showed that if $x^4 \equiv 1 \ (\text{mod } p)$ then
$x^2 \equiv 1 \ (\text{mod } p)$, and since $1 \equiv -1 \ (\text{mod } p)$
only when $p = 2$ we have a contradiction. $\square$
\\

\textit{Solution 2.} \textit{(In this solution do not interpret $(a/b)$
as a fraction. It is instead the Legendre symbol.)}
\vspace{0.1cm}

It suffices to show that $-1$ cannot be a
quadratic residue mod $p$.
Let $a$ be a positive integer and $p$ an odd prime.
Suppose that $a$ and $p$ are coprime (I'm purposefully avoiding the
$a\equiv 0 \ \text{mod } p$ case).
The Legendre symbol is defined like so
$$
\left( \frac{a}{p} \right) =
\begin{cases}
1 & \text{if $a$ is a quadratic residue mod $p$}\\
-1 & \text{if $a$ is not a quadratic residue mod $p$}\\
\end{cases}
$$
Now let $p = 4k+3$ for some $k \in \mathbb{Z}_{\geq0}$. By Euler's
criterion we note that
\[
\begin{aligned}
\left( \frac{p-1}{p} \right) &\equiv (-1)^{\frac{p-1}{2}} \ (\text{mod } p)\\
&\equiv (p-1)^{2k + 1} \ (\text{mod } p)\\
&\equiv -1 \ (\text{mod } p).
\end{aligned}
\]
But if $((p-1) / p)$ were equal to 1 then it wouldn't leave a residue of $-1$
mod $p$ since $1 \equiv -1 \ (\text{mod } p)$ only when $p=2$. Thus it must
be the case that $((p-1) / p) = -1$ and so by definition
$-1$ isn't a quadratic residue mod $p$. $\square$
\\

(c) Let us suppose for the sake of contradiction that $p_1, p_2, \ldots,
p_n$ are all of the primes that are $\equiv 1 \ (\text{mod } 4)$.
Consider the number
$$
P = (2p_1 p_2 \cdots p_n)^2 + 1.
$$
Note that $\gcd(P, p_i) = 1$ due to the Euclidean algorithm.
Due to part (b), $P$ is never divisible by a prime of the form $4k+3$.
But since $P$ is odd it must then have some odd prime factors of the form
$4k+1$ but none of these prime come from the collection $p_1, \ldots, p_n$,
a contradiction. $\blacksquare$
\end{solution}

\end{document}
